\section{CRONOGRAMA DE EXECUÇÃO FÍSICA}
O cronograma divide as etapas entre o planejamento e a qualificação do pré-projeto, no TCC 1, e a construção do MVP ao longo de seis \textit{sprints} quinzenais no TCC 2. Depois disso, o trabalho fecha com o relatório técnico e a defesa (Tabela~\ref{tab:cronograma}).

\begin{table}[!htb]
    \centering
    \caption{Cronograma de Execução Física do EloMaterno}
    \label{tab:cronograma}
    \resizebox{\textwidth}{!}{
    \begin{tabular}{|l|c|c|c|c|c|c|c|c|c|c|c|c|}
        \hline
        \multirow{2}{*}{\textbf{ATIVIDADES}} & \multicolumn{6}{c|}{\textbf{2026 (TCC 1)}} & \multicolumn{6}{c|}{\textbf{2027 (TCC 2)}} \\ \cline{2-13}
        & \textbf{Jul} & \textbf{Ago} & \textbf{Set} & \textbf{Out} & \textbf{Nov} & \textbf{Dez} & \textbf{Jan} & \textbf{Fev} & \textbf{Mar} & \textbf{Abr} & \textbf{Mai} & \textbf{Jun} \\ \hline
        Levantamento bibliográfico e documental & X & X & X & X & X & & & & & & & \\ \hline
        Elaboração e qualificação do pré-projeto & & & X & X & X & & & & & & & \\ \hline
        Refinamento arquitetural e modelagem dos dados & & & & & X & X & & & & & & \\ \hline
        \textit{Sprint 1}: Fundação, Auth JWT e ViaCEP & & & & & & & X & & & & & \\ \hline
        \textit{Sprint 2}: Fórum Comunitário e Filtro PII & & & & & & & X & X & & & & \\ \hline
        \textit{Sprint 3}: Vitrine de Doações e Cloudinary & & & & & & & & X & & & & \\ \hline
        \textit{Sprint 4}: Mural de Pedidos de Ajuda & & & & & & & & & X & & & \\ \hline
        \textit{Sprint 5}: Match Seguro, anonimização e exclusão & & & & & & & & & X & X & & \\ \hline
        \textit{Sprint 6}: Testes ArchUnit, k6 e Deploy & & & & & & & & & & X & & \\ \hline
        Redação final do Relatório Técnico (RTS) & & & & & & & & & & X & X & X \\ \hline
        Revisão com orientador e defesa final & & & & & & & & & & & & X \\ \hline
    \end{tabular}
    }
    \vspace{0.2cm}
    
    {\small Fonte: autor (2026).}
\end{table}

\clearpage